
Analysis of 4,561 models from the Open LLM Leaderboard reveals a dominant latent factor that persists after controlling for model size and training recency. This factor (PC1, $\lambda _1$ = 2.277, 60\% variance explained, p = 0.001) correlates with a behavioral stability proxy ($\rho$ = 0.405) and increases with instruction-tuning (d $\approx$ 2.0).

The cross-benchmark correlation of $\rho$ = 0.80--0.87 across trustworthiness dimensions is not merely a scale artifact. A residual structure persists that instruction-tuning systematically improves. Whether this structure reflects representational stability, general capability spillover, or training data quality cannot be definitively determined from the current evidence.

Key contributions include the first large-scale demonstration of a residual latent factor after confound control (N = 4,561), the proposed BSI-GRC correlation framework, and evidence that instruction-tuning improves both stability proxies and benchmark performance. Critical limitations---synthetic BSI and simulated holdout data---require real behavioral validation and true prospective testing before strong mechanistic conclusions can be drawn.
